% ============================================================
% Chapter 1: Introduction
% CT-MUSIQ Undergraduate Thesis
% ============================================================

\chapter{Introduction}
\label{chap:introduction}

\section{Background and Motivation}
\label{sec:intro_background}

Computed Tomography (CT) has established itself as one of the most indispensable modalities
in modern clinical medicine. Since its clinical introduction by Hounsfield and Cormack in
the early 1970s, CT imaging has undergone decades of technological evolution, culminating
in today's multi-detector row scanners capable of acquiring full-body volumetric datasets
in under ten seconds \cite{hounsfield1973}. The exceptional anatomical resolution, rapid
acquisition time, and wide clinical accessibility of CT have made it an irreplaceable tool
for emergency triage, oncological staging, neurological assessment, and pre-surgical
planning. Globally, hundreds of millions of CT scans are performed annually, generating an
enormous volume of cross-sectional healthcare imaging data.

Despite its clinical indispensability, CT imaging carries an inherent risk that has long
concerned the medical physics and radiology communities: exposure to ionising X-ray
radiation. Unlike natural background radiation, medical X-ray exposure is controllable,
and the clinical community has adopted the ALARA principle (As Low As Reasonably
Achievable) as the governing framework for dose management \cite{alara_ref}. Under this
principle, Low-Dose CT (LDCT) protocols have been aggressively developed and deployed,
achieving dose reductions of 50--80\% compared to conventional normal-dose CT (NDCT)
without compromising diagnostic yield for a broad range of indications. The clinical
validation of LDCT has been particularly compelling in the context of lung cancer screening
programmes, where multi-institutional observational studies have demonstrated that
ultra-low-dose protocols can maintain the sensitivity necessary for early-stage nodule
detection \cite{ldct_lung_screening}.

However, dose reduction is not without consequence. The statistical nature of X-ray
detection imposes a fundamental physical trade-off: reducing the number of X-ray photons
incident on the detector array directly degrades the signal-to-noise ratio of the acquired
sinogram. Specifically, the photon count at each detector cell follows a Poisson
distribution, and the variance of this distribution scales linearly with the mean photon
count. When the mean count is reduced, the relative noise (the coefficient of variation)
increases as the inverse square root of the count, producing a complex, spatially
non-stationary noise field in the reconstructed image. Superimposed on the quantum noise
is the Gaussian-distributed electronic readout noise of the data acquisition system,
creating the well-known ``Poisson-plus-Gaussian'' degradation model of LDCT \cite{zhang2024review}.
In the image domain, this manifests as conspicuous streak artefacts, grainy noise texture,
loss of low-contrast lesion conspicuity, and blurring of subtle anatomical boundaries.

The clinical impact of these degradations is not uniform across tissues or disease states.
In brain CT, the primary diagnostic challenge is to differentiate structures with inherently
small HU differences: grey matter ($\sim$40~HU), white matter ($\sim$35~HU), cerebrospinal
fluid (0--15~HU), and pathological processes such as hyperacute haemorrhage (50--80~HU),
oedema, and early infarction. The narrow Hounsfield unit dynamic range of brain tissue
renders it exquisitely sensitive to any noise perturbation that narrows the effective
contrast-to-noise ratio, making brain LDCT quality assessment a particularly demanding
and clinically high-stakes problem \cite{brainCT_window}.

Conventionally, CT image quality is evaluated by trained radiologists who inspect each
acquired scan for diagnostic acceptability, using structured criteria covering noise
level, artefact severity, contrast resolution, and overall lesion visibility. This gold
standard is both the most clinically valid and the most operationally burdensome approach:
it is labour-intensive, subject to significant inter-observer and intra-observer
variability, and fundamentally non-scalable to the throughput demands of modern
high-volume CT facilities or retrospective research analysis pipelines. As deep learning
post-processing algorithms increasingly generate large numbers of reconstructed images
for research evaluation or prospective quality control, the manual radiological audit
bottleneck becomes acutely limiting \cite{herath2025}.

This gap between clinical need and operational feasibility provides the primary motivation
for \textbf{automated, no-reference Image Quality Assessment (NR-IQA)} for LDCT. An
effective NR-IQA system would accept a single CT slice as input and predict a
perceptual quality score aligned with radiologist assessment, enabling:

\begin{itemize}
    \item Real-time scan rejection or flagging in the scanner room, prompting immediate
          re-acquisition at higher dose if a scan is below diagnostic threshold;
    \item Large-scale retrospective quality auditing of CT archives without radiologist
          involvement;
    \item Objective, reproducible benchmarking of competing LDCT denoising and
          reconstruction algorithms;
    \item Integration into clinical Picture Archiving and Communication Systems (PACS)
          as an automated quality gate before image distribution to radiologists.
\end{itemize}

Achieving these goals requires overcoming a series of fundamental technical challenges,
which have resisted satisfactory resolution by prior computational approaches, as reviewed
in detail in Chapter~\ref{chap:literature_review}.

\section{Problem Statement}
\label{sec:intro_problem}

The formal task addressed in this thesis is defined as follows. Given a single axial brain
CT slice $\mathbf{x} \in \mathbb{R}^{H \times W}$ (with $H = W = 512$) acquired under
low-dose protocol, the objective is to learn a function:

\begin{equation}
    f_\theta : \mathbb{R}^{H \times W} \rightarrow [0, 4]
\end{equation}

that maps the input CT image to a continuous quality score $\hat{q} = f_\theta(\mathbf{x})$
maximally correlated with the consensus perceptual quality score $q \in [0, 4]$ assigned
by a panel of expert radiologists using a five-point Likert scale. Crucially, $f_\theta$
operates in the \textbf{no-reference} (or blind) setting: no paired high-quality reference
image is available either during training inference, reflecting the clinical reality that
acquiring such reference images would itself require an additional radiation dose.

The specific challenges that distinguish this problem from standard NR-IQA on natural
photographs include:

\begin{enumerate}
    \item \textbf{Domain specificity of degradations.} CT image noise arises from physical
    quantum mechanics and detector electronics rather than from digital compression or
    optical blur. The resulting non-stationary Poisson-driven noise does not conform to
    the Natural Scene Statistics (NSS) assumptions underlying classical NR-IQA algorithms
    \cite{brisque_original, niqe_original}.

    \item \textbf{Calibrated Hounsfield Unit scale.} Unlike RGB natural images, CT pixel
    values are physical quantities with precise clinical meaning. The clinically relevant
    intensity range for brain tissue spans only 80~HU (0 to 80~HU after brain soft-tissue
    windowing), meaning that the specific preprocessing applied to the image before
    assessment fundamentally shapes the perceptual quality distribution seen by both the
    human radiologist and any automated model.

    \item \textbf{Global artefact structure.} LDCT streak artefacts arising from photon
    starvation propagate across the entire image cross-section along Radon transform
    directions, creating long-range spatial correlations that local convolutional kernels
    cannot model with finite receptive fields.

    \item \textbf{Extremely limited annotated training data.} Radiologist annotation of
    CT images is expensive and ethically constrained. The largest publicly available
    annotated LDCT IQA dataset, released by the LDCTIQAC 2023 challenge, contains only
    1,000 labelled images. Training deep neural networks of the scale required for
    state-of-the-art IQA on such limited data demands rigorous regularisation strategies.

    \item \textbf{Input resolution paradox.} High-resolution CT images (512~$\times$~512
    pixels) must be processed to capture fine-grained noise signatures, yet standard
    Vision Transformer architectures require fixed-resolution inputs and cannot natively
    process variable-length patch sequences across multiple scales simultaneously.
\end{enumerate}

\section{Proposed Approach: CT-MUSIQ}
\label{sec:intro_approach}

This thesis proposes \textbf{CT-MUSIQ}, a purpose-built NR-IQA framework that addresses
each of the above challenges through a cohesive set of architectural and training
innovations. CT-MUSIQ is constructed as a principled medical domain adaptation of the
Multi-Scale Image Quality Transformer (MUSIQ) \cite{ke2021musiq}, an architecture
originally designed for natural image quality assessment that was selected as the
foundation for this work due to its unique ability to process images at multiple
resolutions without interpolation-induced information loss.

The core design philosophy of CT-MUSIQ rests on three guiding principles:

\textbf{Principle 1: Clinical fidelity through domain-specific preprocessing.} Rather
than applying standard ImageNet normalisation to CT images, CT-MUSIQ enforces the brain
soft-tissue CT window transformation (level 40~HU, width 80~HU) that replicates the
exact intensity range presented to annotating radiologists during the labelling process.
This ensures that the model learns quality representations in the same perceptual space
as the human ground truth labels.

\textbf{Principle 2: Multi-scale global receptive field from the first layer.} By
extracting non-overlapping $32 \times 32$ patches from the CT image at both $224 \times
224$ and $384 \times 384$ resolutions, CT-MUSIQ simultaneously provides the Transformer
encoder with coarse-scale global structural context (49 patches per image) and fine-scale
local noise texture (144 patches per image), without any spatial downsampling that would
destroy frequency information critical for noise assessment.

\textbf{Principle 3: Regularisation through scale consistency.} Given the small
training set, the model must be prevented from learning scale-specific overfitting
artefacts. The scale-consistency KL divergence loss, a novel contribution of this thesis,
penalises disagreement between the quality assessments of individual scale-specific
prediction heads and the global aggregated prediction, enforcing that each scale learns
a coherent, internally consistent quality representation.

\section{Contributions}
\label{sec:intro_contributions}

The principal scientific and engineering contributions of this thesis are as follows:

\begin{enumerate}
    \item \textbf{CT-MUSIQ architecture.} A complete, end-to-end trainable NR-IQA
    architecture adapted from MUSIQ for brain CT imaging, incorporating multi-scale patch
    tokenisation, hash-based positional encoding, and per-scale prediction heads. The
    architecture is designed to operate within 6~GB of GPU VRAM on consumer hardware
    using mixed-precision training.

    \item \textbf{Brain CT windowing adaptation.} A principled methodological argument
    for and implementation of brain soft-tissue windowing as a preprocessing step that
    bridges the perceptual space between radiologist annotation and model training,
    constituting a significant departure from the abdominal windowing used in the original
    LDCTIQAC challenge.

    \item \textbf{Scale-consistency KL regularisation.} A novel composite loss function
    that combines mean squared error regression, pairwise margin ranking, and
    Kullback--Leibler divergence between per-scale and global quality prediction
    distributions, verified through systematic ablation studies to provide measurable
    improvement over MSE-only training on the small LDCT dataset.

    \item \textbf{Two-stage fine-tuning protocol.} A carefully engineered training
    strategy that first warms up domain-specific components with the Transformer encoder
    frozen, then jointly fine-tunes all parameters with cosine-annealed learning rate
    decay and exponential moving average weight smoothing, achieving stable convergence
    on 700 training samples.

    \item \textbf{Comprehensive empirical evaluation.} A full benchmark of CT-MUSIQ
    against both published LDCTIQAC 2023 leaderboard methods and a purpose-built strong
    baseline (Swin-T + ResNet-50 ensemble), with ablation studies confirming the
    individual contribution of each proposed component, and an analysis of the
    complementary strengths of the two approaches.

    \item \textbf{Open-source reproducible implementation.} A complete Python/PyTorch
    implementation of CT-MUSIQ, including all training, evaluation, and ablation scripts,
    centralised configuration management, and deterministic reproducibility through
    fixed random seeds, released as a fully documented open-source repository.
\end{enumerate}

\section{Thesis Structure}
\label{sec:intro_structure}

The remainder of this thesis is organised as follows.

\textbf{Chapter~\ref{chap:literature_review}: Literature Review} surveys the relevant
prior work across five thematic domains: the clinical evolution of CT reconstruction
algorithms; the theoretical and practical limitations of full-reference IQA in the
medical imaging context; the failure modes of classical NSS-based NR-IQA when applied
to CT data; the development of deep learning-based NR-IQA methods including CNN-based and
Transformer-based architectures; and the emergence of vision-language models as the
current frontier of IQA research. The chapter concludes with a precise formulation of the
research gap addressed by this thesis.

\textbf{Chapter~\ref{chap:methodology}: Methodology} provides a comprehensive technical
description of all components of CT-MUSIQ: the dataset and preprocessing pipeline,
the model architecture in full mathematical detail, the composite loss function, the
two-stage training procedure, data augmentation strategy, evaluation protocol, and the
comparison baseline architecture.

\textbf{Chapter~\ref{chap:experiments}: Experiments} describes the experimental design
including implementation details, hardware specifications, training hyperparameters, and
the complete ablation study protocol. The design of each ablation configuration is
motivated by the specific architectural hypotheses it tests.

\textbf{Chapter~\ref{chap:results}: Results and Discussion} presents the quantitative
performance of CT-MUSIQ on the LDCTIQAC 2023 test set relative to all published
leaderboard methods and the Agaldran Combo baseline, followed by detailed analysis of
the ablation study results, examination of correlation metric components, and discussion
of the implications of the findings for clinical deployment.

\textbf{Chapter~\ref{chap:conclusion}: Conclusion and Future Work} synthesises the key
contributions and findings of this thesis, critically examines the limitations of the
proposed approach, and outlines a roadmap of promising directions for future research
in automated CT quality assessment.

\textbf{Appendices} provide supplementary material including the complete hyperparameter
summary table, training curve data, algorithm pseudocode, and implementation notes for
reproducing all experiments.

\section{Scope and Limitations}
\label{sec:intro_scope}

This thesis focuses specifically on axial brain CT slices from the LDCTIQAC 2023 dataset
and does not claim generalisation to other CT modalities (e.g., chest, abdominal, or
orthopaedic CT), multi-slice volumetric assessment, or other medical imaging modalities.
The evaluation is conducted exclusively on a single publicly available benchmark dataset;
multi-institutional validation across different scanner manufacturers, acquisition
protocols, and patient populations represents an important direction for future work.

The CT images used in this study are pre-processed and normalised prior to release by the
challenge organisers; direct application to raw DICOM files from clinical PACS systems
would require additional preprocessing steps including HU calibration verification and
patient table removal.

Furthermore, the Agaldran Combo baseline used in this thesis as a within-paper comparison
is not the original ``agaldran'' team's submission to the LDCTIQAC 2023 challenge, but
rather an independent reimplementation motivated by their published architectural choices,
trained under the same conditions as CT-MUSIQ. Direct comparison with the original
challenge submission is therefore approximate, though the reported leaderboard numbers
are cited for reference throughout.


\section{Clinical Context and Significance}
\label{sec:intro_clinical}

To fully appreciate the significance of the automated IQA problem addressed in this
thesis, it is important to understand the clinical workflow in which CT image quality
assessment occurs and the specific consequences of quality failures in routine clinical
practice.

\subsection{The CT Acquisition Workflow}

In a typical clinical CT examination, the following sequence of events occurs: (1) a
clinician orders a CT scan for a patient based on clinical indications; (2) the patient
is positioned on the CT scanner table and the acquisition parameters (kVp, mAs, scan
duration) are set by the radiographer according to local protocol, often using
dose-modulation algorithms; (3) the scanner acquires the raw sinogram data and performs
on-line reconstruction using the configured algorithm (FBP, IR, or DLIR); (4) the
reconstructed image series is transferred to the hospital PACS; (5) a radiologist accesses
the images from PACS, applies windowing, and performs a diagnostic assessment;
(6) a radiological report is generated and communicated to the referring clinician.

In this workflow, quality assessment traditionally occurs implicitly at step~(5): the
radiologist observes the image quality as part of the diagnostic reading and may note
quality issues in the report or, in severe cases, request a repeat scan. However, by
step~(5), the patient has already left the scanner room, and a repeat scan requires
re-booking the patient, additional radiation exposure, additional scanner time, and
delay in diagnosis. An automated quality gate at step~(3) or (4)---before the patient
leaves the room---would enable immediate repeat acquisition when necessary, avoiding
all of these downstream costs.

\subsection{The Economic and Radiation Safety Case}

The economic case for automated IQA is straightforward. A missed poor-quality scan that
reaches the radiologist requires: (a) radiologist time to assess and report on a
subdiagnostic image; (b) communication with the referring clinician to explain the
limitation; (c) in many cases, a repeat scan booking, which incurs scanner time,
staff time, and patient transport costs; and (d) a delayed diagnosis that may have
clinical consequences if the pathology was urgent. The cost of a single repeat CT
examination varies by region and modality but typically ranges from \$200--\$1,000.
Even a 1\% reduction in repeat scan rates for a facility performing 50,000 CT scans
per year would save 500 repeat scans annually---a significant operational benefit.

From a radiation safety perspective, eliminating unnecessary repeat scans directly
reduces the aggregate population radiation dose. Modern CT scanner dose modulation
algorithms already minimise dose for any given diagnostic target, but a missed
poor-quality scan followed by a repeat effectively doubles the patient's radiation
exposure for that examination. An IQA system that prevents this scenario contributes
directly to the ALARA principle that governs clinical radiology practice.

\subsection{Research Benchmarking Applications}

Beyond clinical quality control, automated IQA is essential for the reproducible
benchmarking of competing CT reconstruction and denoising algorithms. The literature
review (Chapter~\ref{chap:literature_review}) documents the rapidly expanding landscape
of DLIR methods, each claiming performance improvements over prior art. Without an
objective, perceptual quality metric aligned with radiologist judgement, comparisons
between methods rely either on subjective visual inspection (which is non-reproducible
and expert-dependent) or on non-perceptual metrics such as PSNR and SSIM (which often
do not correlate with diagnostic quality). CT-MUSIQ provides a reproducible, scalable,
radiologist-aligned quality assessment tool that can serve as an objective benchmark
for evaluating LDCT reconstruction quality across large image databases.
